\documentclass[a4paper,10pt]{article}
\usepackage[UTF8,fontset=fandol]{ctex}
\usepackage{geometry}
\usepackage{titlesec}
\usepackage{enumitem}
\usepackage{hyperref}
\usepackage{xcolor}
\geometry{a4paper, top=1.2cm, bottom=1.2cm, left=1.5cm, right=1.5cm}
\pagestyle{empty}
\setlength{\parindent}{0pt}
\linespread{1.08}
\titleformat{\section}{\large\bfseries}{}{0em}{}[\titlerule]
\titlespacing*{\section}{0pt}{0.4cm}{0.2cm}
\begin{document}
\begin{center}
{\Huge \textbf{雷军（教学演示版）}} \\
\vspace{0.15cm}
\textbf{意向岗位：京东采销} \\
\vspace{0.1cm}
电话：138-0000-0000 \quad | \quad 邮箱：demo@example.com \quad | \quad 地点：武汉
\vspace{0.08cm}
\textbf{岗位关键词：} 复盘优化 / 跨部门协作 / 数据分析 / 项目管理 / 风险管理 / 商业分析 / 活动运营 / 行业研究
\end{center}
\section*{教育背景}
\textbf{武汉大学} \hfill 1987.09 - 1991.06 \\
计算机科学与工程 \hfill 本科（应届） \\
\textbf{核心亮点：} 公开报道提及：2年修完本科全部学分；大学阶段形成技术实现与商业落地并重的实践导向
\section*{实习/项目经历}
\textbf{BITLOK 加密软件（校园商业化项目）} \hfill 1989.03 - 1990.12 \\
\textit{联合开发者} \hfill 武汉
\begin{itemize}[leftmargin=*, itemsep=2pt, parsep=0pt, topsep=2pt]
    \item \textbf{协同推进}：在有限资源下保证交付质量；拆分任务并进行每日进度同步，确保关键模块按期交付；高压环境下保持稳定产出；关键指标：\textbf{连续高强度推进}、\textbf{关键节点无延期}。
    \item \textbf{项目统筹}：完成加密软件首版交付并验证商业可行性；与伙伴连续15天高强度开发与测试，搭建安装、加密、恢复的完整流程；项目实现从技术实现到商业化落地；关键指标：\textbf{15天完成首版开发}、\textbf{单套定价2000元以上}。
    \item \textbf{复盘优化}：根据反馈快速迭代核心功能；围绕安装成本、易用性和稳定性持续优化，形成版本迭代机制；产品形成长期销售能力；关键指标：\textbf{持续销售约6-7年}。
\end{itemize}
\textbf{校园汉卡项目（失败复盘案例）} \hfill 1990.01 - 1990.10 \\
\textit{项目负责人} \hfill 武汉
\begin{itemize}[leftmargin=*, itemsep=2pt, parsep=0pt, topsep=2pt]
    \item \textbf{复盘优化}：快速止损并输出复盘结论；重新梳理需求验证、现金流和渠道策略，形成复盘文档；沉淀可复用的风险管理方法；关键指标：\textbf{完成一次完整失败复盘}、\textbf{明确3条运营底线}。
    \item \textbf{项目统筹}：组织打样与首批备货并跑通销售；以借款启动项目并推进生产、库存和渠道动作；建立了从供给端到销售端的完整运营视角；关键指标：\textbf{借款14000元启动}、\textbf{首批库存2000套}。
\end{itemize}
\textbf{高密度学习与自驱成长（大学阶段）} \hfill 1987.09 - 1989.09 \\
\textit{学习项目执行者} \hfill 武汉
\begin{itemize}[leftmargin=*, itemsep=2pt, parsep=0pt, topsep=2pt]
    \item \textbf{复盘优化}：建立科技创业导向的能力地图；通过系统阅读与实践，将长期目标拆解为阶段性任务；形成长期主义与结果导向的工作方法；关键指标：\textbf{持续多年自驱实践}。
    \item \textbf{项目统筹}：压缩学习周期并提升产出效率；制定高密度学习计划并在图书馆与机房长期执行；提前完成学业要求并转入实战阶段；关键指标：\textbf{2年修完4年学分}。
\end{itemize}
\section*{项目补充}
\textbf{校园软件商业化实战（BITLOK）} \hfill 1989.03 - 1990.12 \\
\begin{itemize}[leftmargin=*, itemsep=2pt, parsep=0pt, topsep=2pt]
    \item \textbf{项目落地}：将技术项目推进到可收费产品形态，验证需求到交易的闭环路径。
    \item \textbf{项目落地}：在有限资源下持续迭代并保持长期销售能力。
\end{itemize}
\section*{技能与工具}
\begin{itemize}[leftmargin=*, itemsep=2pt, parsep=0pt, topsep=2pt]
    \item \textbf{技术基础：} C语言、基础汇编、早期PC软件开发流程
    \item \textbf{运营能力：} 需求洞察、项目推进、复盘改进、商业化思维
    \item \textbf{个人特质：} 强执行力、抗压推进、结果导向
\end{itemize}
\section*{自我评价}
具备强执行力与结果导向，能够在高节奏环境中推进复杂任务落地；重视数据与复盘，能持续优化策略并稳定交付业务结果。
\end{document}
